\clearpage
\section{Solutions to the exercises}
\label{sec-solutions-7-representations}

These solutions use invariant subspaces, averaging, and character identities
to analyze representations and their applications. Particular attention is
given to the hypotheses behind decomposition, the construction of polynomial
invariants, and the normalization of finite-field wavelet reconstruction.

\begin{correction}{exo-reductibility-decomposability}
The line $Ke_1$ is invariant. It is the only invariant line: if
$v=ae_1+be_2$ with $b\ne0$, invariance under $\rho(1)$ would force
$\rho(1)v=v$, by comparing second coordinates, whereas their difference
is $be_1\ne0$. Thus no direct sum of two invariant lines is possible.
For $K=\CC$, the acting additive group is infinite. For finite $K$ of
characteristic $p$, its order is a power of $p$, so division by $|G|$
in Maschke's averaging argument is impossible. This is a characteristic
obstruction, not an algebraic-closure obstruction.
\end{correction}

\begin{correction}{exo-stationary-operators}
Set $\varphi=A\delta_1$. Since $\delta_g=\tau_g\delta_1$,
\[
 Af=\sum_g f(g)A\delta_g=\sum_g f(g)\tau_g\varphi=f*\varphi.
\]
Evaluating this identity at $f=\delta_1$ proves uniqueness. Right
convolution commutes with the left regular action by associativity.
Replacing it by left convolution would require additional commutation
conditions on the kernel in a nonabelian group.
\end{correction}

\begin{correction}{exo-irred-representation-degre-1}
The twisted matrices satisfy the representation law because $\eta$ is
a scalar homomorphism. A subspace is invariant under them exactly when
it is invariant under $\rho$, since $\eta(g)\ne0$. The twist is therefore
irreducible, with character $\eta\chi_\rho$. Equivalently,
$|\eta(g)|=1$ makes the character norm unchanged.
For $G=\mathfrak S_3$, the standard character has values $(2,0,-1)$
on the identity, transpositions, and three-cycles. Multiplying by sign
leaves it unchanged, so the twist is isomorphic to the same representation.
If $\rho$ is trivial, its twist is $\eta$ itself. Thus nontriviality of
$\eta$ alone does not ensure a new representation.
\end{correction}

\begin{correction}{exo-representation-product}
Let $\rho_i$ and $\sigma_j$ be the irreducibles of $G$ and $H$.
The external tensor products
\[
 (g,h)\longmapsto\rho_i(g)\otimes\sigma_j(h)
\]
have characters $\chi_i(g)\eta_j(h)$. The normalized inner product
on the product group factors, so these characters are orthonormal.
They therefore give pairwise nonisomorphic irreducibles. Finally,
\[
 \sum_{i,j}(\dim V_i\dim W_j)^2
       =\left(\sum_i\dim V_i^2\right)\left(\sum_j\dim W_j^2\right)
       =|G||H|.
\]
The degree-squares identity proves completeness.
\end{correction}

\begin{correction}{exo-action-polynomials}
\begin{enumerate}
\item Permutation matrices identify the action with permutation of
indeterminates. The invariant ring is generated freely by the elementary
symmetric polynomials
$\sigma_j=\sum_{i_1<\cdots<i_j}X_{i_1}\cdots X_{i_j}$, $1\leq j\leq n$.
\item Independent sign changes force both exponents to be even, giving
$K[X,Y]^{V_4}=K[X^2,Y^2]$. The two generators are algebraically independent.
For $C_2$, only the total degree must be even, so
$K[X,Y]^{C_2}=K[X^2,XY,Y^2]$. These three degree-two elements are needed,
but expressions in them are not unique because $(XY)^2=X^2Y^2$.
\item If $P$ is invariant and $P=\sum_\alpha c_\alpha X^\alpha$, then
$P=R_G(P)=\sum_\alpha c_\alpha R_G(X^\alpha)$.
\item Expand each $U_A^k$ by the multinomial theorem and sum over $A$.
This gives the asserted identity for $S_k$. For $N$ variables, Newton's
identities in characteristic zero express every symmetric polynomial
in the first $N$ power sums. Applied to the $N$ quantities $(U_A)$,
they give $S_k=F_k(S_1,\ldots,S_N)$ with coefficients in $K$.
Compare the coefficient of $u^\alpha$ for $|\alpha|=k$: the left side
is $N\binom{k}{\alpha}R_G(X^\alpha)$; the right side is a polynomial
in averages of degrees at most $N$. Division by the nonzero integer
$N\binom{k}{\alpha}$ proves the bound. Characteristic zero justifies
both divisions in Newton's identities and this final division.
\item Odd-degree averages vanish because $-I\in C_4$. At degree two,
$R(X^2)=R(Y^2)=(X^2+Y^2)/2$ and $R(XY)=0$. At degree four,
\[
 \begin{aligned}
 R(X^4)=R(Y^4)&=(X^4+Y^4)/2,\\
 R(X^3Y)&=(X^3Y-XY^3)/2,\\
 R(X^2Y^2)&=X^2Y^2,\qquad R(XY^3)=-R(X^3Y).
 \end{aligned}
\]
Noether's bound therefore gives generators
\[
 P=X^2+Y^2,\quad Q=X^2Y^2,\quad R=XY(X^2-Y^2),
\]
since $X^4+Y^4=P^2-2Q$. They satisfy $R^2=Q(P^2-4Q)$.
One degree-two generator and two additional degree-four generators are
necessary: the invariant spaces at these degrees have dimensions one
and three, respectively, whereas $P$ alone produces only $P^2$ at degree four.
\end{enumerate}
\end{correction}

\begin{correction}{exo-thm-molien}
\begin{enumerate}
\item Linear substitutions preserve every homogeneous degree. Equality
of two polynomials is equivalent to equality of their homogeneous
components, which proves the claim.
\item The Reynolds operator on $V_s$ projects onto $V_s^G$, hence its
trace is $a_s$. With the specified inverse matrix and monomial basis,
\[
 A^{[2]}=\begin{pmatrix}
 a^2&ac&c^2\\2ab&ad+bc&2cd\\b^2&bd&d^2
 \end{pmatrix}.
\]
These entries come from $(aX+bY)^2$, $(aX+bY)(cX+dY)$, and
$(cX+dY)^2$. Using entries of $A$ in place of $A^{-1}$ would describe
a different substitution convention.
\item A finite-order complex matrix is diagonalizable. If
$\mu_1,\ldots,\mu_n$ are the eigenvalues of $A^{-1}$, those on degree
$s$ are $\mu^\alpha$ for $|\alpha|=s$. Therefore
\[
 \sum_{s\geq0}\operatorname{tr}(A^{[s]})t^s
   =\prod_{j=1}^n(1-t\mu_j)^{-1}=\det(I-tA^{-1})^{-1}.
\]
Average over the group and replace $A$ by $A^{-1}$ to obtain Molien's
formula, including its factor $1/|G|$.
\item For the reflection group $G_1$ of order two,
\[
 \Phi_1(t)=\frac12\left(\frac1{(1-t)^2}+\frac1{1-t^2}\right)
          =\frac1{(1-t)(1-t^2)}.
\]
The group $G_2$ in the self-dual-code exercise includes the even-weight
reflection $\operatorname{diag}(1,-1)$ and has order sixteen. Its eight
rotations and eight reflections give
\[
 \Phi_2(t)=\frac1{16}\left(
    \sum_{j=0}^7\frac1{1-2t\cos(j\pi/4)+t^2}
       +\frac8{1-t^2}\right)
       =\frac1{(1-t^2)(1-t^8)}.
\]
After verifying that proposed invariants are algebraically independent,
the subalgebra they generate has a known homogeneous dimension in each
degree. Equality with the Molien coefficients then proves generation.
A denominator resembling that of a polynomial ring is not by itself
such a proof; explicit invariants and their independence are required.
\end{enumerate}
\end{correction}

\begin{correction}{exo-lemma-cauchy-frobenius}
\begin{enumerate}
\item The trace of a permutation matrix is its number of fixed points,
so $\dotp{\chi_\pi}{1}=|G|^{-1}\sum_g|X_g|$.
\item The invariant functions are constant on each orbit, and the orbit
indicators form a basis. Thus $\dim V^G$ equals the number of orbits;
the Reynolds trace gives the stated average. Alternatively,
\[
 \sum_g|X_g|=\#\{(g,x):gx=x\}=\sum_x|G_x|
           =\sum_{\text{orbits }O}|O|\frac{|G|}{|O|}.
\]
\item A rotation by $2\pi j/6$ has $\gcd(6,j)$ cycles and fixes
$2^{\gcd(6,j)}$ binary colorings. The six rotation counts are
$64,2,4,8,4,2$. Three reflections through opposite vertices fix
$2^4=16$ colorings each; three through opposite edge midpoints fix
$2^3=8$ each. Hence the number of necklaces is
\[
 \frac{64+2\cdot2+2\cdot4+8+3\cdot16+3\cdot8}{12}=13.
\]
The half-turn contribution $8$ must be included, and adjacent rotations
are separated by $\pi/3$, not $\pi/6$.
\item A matrix commutes with the permutation action precisely when its
entries are constant on the orbits of the diagonal action on $X\times X$.
Double transitivity gives exactly two orbits, the diagonal and its
complement, hence $\dim\End_G(V)=2$ and
$\dotp{\chi_\pi}{\chi_\pi}=2$. Transitivity gives $\dim V^G=1$.
Write $V=\CC\mathbf1\oplus W$ with no invariants in $W$; character
orthogonality then gives $\norm{\chi_W}_2^2=2-1=1$.
Thus $W$ is irreducible. The natural action of $\mathfrak S_n$ is
doubly transitive for $n\geq2$.
\end{enumerate}
\end{correction}

\begin{correction}{exo-representation-theory-numbers}
\begin{enumerate}
\item Conjugating the class sum permutes its terms, so it is an
intertwiner. Schur's lemma makes it scalar; taking its trace yields
$r_C=|C|\chi_\rho(C)/d$.
\item Character orthogonality gives
\[
 |G|=\sum_C|C|\chi_\rho(C)\chi_\rho(C^{-1})
      =d\sum_Cr_C\chi_\rho(C^{-1}).
\]
\item If $x'=hxh^{-1}$, simultaneous conjugation
$(a,b)\mapsto(hah^{-1},hbh^{-1})$ bijects the factorizations of $x$
with those of $x'$. The number is therefore constant on the class.
\item Multiply two class sums and collect products by their conjugacy
class. Applying $\rho$ gives $r_Cr_D=\sum_Ea_{CD}^Er_E$.
Their integer span is consequently closed under multiplication and
contains $r_{\{1\}}=1$. It is finitely generated as a $\ZZ$-module,
not just as a ring.
\item This module is torsion-free inside $\CC$, hence free of finite
rank. Multiplication by $r_C$ has an integer matrix in a $\ZZ$-basis.
Its monic characteristic polynomial annihilates the operator and, upon
application to one, annihilates $r_C$. Thus $r_C$ is an algebraic integer.
Each $\rho(g)$ has finite order, so its eigenvalues are roots of unity;
its trace is a \emph{sum} of such roots and is an algebraic integer.
The formula in part~2 makes $|G|/d$ an algebraic integer. Since it is
rational, it is an integer, proving divisibility.
\end{enumerate}
\end{correction}

\begin{correction}{exo-det-groupe}
\begin{enumerate}
\item Acting on $\delta_h$, the element $X=\sum_aX_a\delta_a$ gives
$\sum_aX_a\delta_{ah}$. Its coordinate at $g$ is $X_{gh^{-1}}$, so
this is exactly the regular-action matrix. With the normalized inner
product, that coordinate is $|G|\dotp{X*\delta_h}{\delta_g}$, not the
inner product alone.
\item Direct sums give block-diagonal matrices, whose determinants multiply.
\item The map $\rho_i$ is the composition of the Fourier isomorphism
with projection onto its $i$th full matrix block, hence is surjective.
A linear relation among its entry forms would then vanish on every
matrix unit, forcing all coefficients to be zero.
\item Induct on $n$, the case $n=1$ being immediate. Expand along the
last row, regarding its entries as indeterminates over the polynomial
ring $R$ in the first $n-1$ rows. Each cofactor is a generic determinant
of size $n-1$, irreducible by induction, and two cofactors obtained by
deleting different columns are nonassociate. Their greatest common divisor
in $R$ is therefore one. The row expansion is a primitive polynomial
over this UFD. Over the fraction field of $R$, it is a nonconstant
linear polynomial and hence irreducible. Gauss's lemma concludes.
Primitivity is essential: a polynomial linear in a variable may still
have a nonconstant common factor in its coefficients.
\item Complete the $d_i^2$ independent entry forms to a basis of all
linear forms in $(X_g)$. In the resulting invertible linear coordinates,
$\Theta_{\rho_i}$ is the generic determinant in the first $d_i^2$
variables, independent of the remaining variables. Irreducibility is
preserved by adjoining these variables and by the invertible substitution.
\item Since $\rho_i(1)=I$, write $\rho_i(X)=X_1I+B$, with
$B=\sum_{g\ne1}X_g\rho_i(g)$. The coefficient of $X_1^{d_i-1}$ in
$\det(X_1I+B)$ is $\operatorname{tr}B=\sum_{g\ne1}\chi_i(g)X_g$.
The degree gives $d_i=\chi_i(1)$, so the polynomial determines the
entire character, hence the isomorphism class. Proportional such factors
must have the same degree and, because they are monic in $X_1$, must
be equal. Therefore distinct irreducible classes give nonproportional factors.
\item The regular representation contains $V_i$ with multiplicity $d_i$.
Apply the direct-sum determinant formula and the preceding irreducibility
and distinctness results.
\end{enumerate}
\end{correction}

\begin{correction}{exo-group-affine-finite-body}
\begin{enumerate}
\item Composition gives
\[
 (b,a)(d,c)=(b+ad,ac),\quad 1=(0,1),\quad
 (b,a)^{-1}=(-a^{-1}b,a^{-1}).
\]
This is the semidirect product with $\FF_p^*$ acting on the additive
group by multiplication. Its order is $p(p-1)$.
\item The action on points is $x\mapsto ax+b$, and the action on
functions is its inverse pullback. Thus it respects composition and
permutes coordinates, preserving the normalized Hermitian product.
\item The mean is unchanged by any coordinate permutation, so $E$ is
invariant. Write $e_k(x)=\omega_p^{kx}$, with
$\omega_p=e^{2\pi\imath/p}$. Then
\[
 \pi(b,a)e_k=\omega_p^{-ka^{-1}b}e_{ka^{-1}}.
\]
The nonzero $k$ form a basis of $E$. If $a\ne1$, no such frequency is
fixed, so the character is zero. For $a=1$, its value is $p-1$ when
$b=0$ and $-1$ when $b\ne0$. Hence its squared character norm is
\[
 \frac{(p-1)^2+(p-1)}{p(p-1)}=1,
\]
proving irreducibility. Alternatively, an affine map is uniquely determined
by its values at two distinct points, so the action is doubly transitive
and the preceding counting exercise applies.
\end{enumerate}
\end{correction}

\begin{correction}{exo-wavelets-finite-fields}
\begin{enumerate}
\item A change of variable gives
$\wh\psi_{b,a}(k)=\omega_p^{-bk}\wh\psi(ak)$. Plancherel then yields
\[
 \Ww_\psi f(b,a)=\frac1p\sum_{k\in\FF_p}
              \wh f(k)\overline{\wh\psi(ak)}\omega_p^{bk}.
\]
For each fixed dilation, this is an inverse ordinary DFT.
\item Define the unnormalized synthesis
$Sy=\sum_{a\ne0,b}y(b,a)\psi_{b,a}$. Taking its Fourier transform and
summing over translations gives
\[
 \widehat{S\Ww_\psi f}(k)=\wh f(k)\sum_{a\ne0}|\wh\psi(ak)|^2.
\]
For $k\ne0$, multiplication by $k$ permutes $\FF_p^*$, so the last
sum is $\sum_{u\ne0}|\wh\psi(u)|^2=c_\psi$ when $\psi$ has mean zero.
Parseval gives $c_\psi=p\sum_x|\psi(x)|^2>0$. At $k=0$ the result
vanishes. Thus $S\Ww_\psi f=c_\psi f$ on $E$, and for arbitrary $f$
the reconstruction is $f-p^{-1}\sum_xf(x)$.
\item For a general wavelet, the synthesis-analysis multiplier equals
$(p-1)|\wh\psi(0)|^2$ at zero and $\sum_{u\ne0}|\wh\psi(u)|^2$ elsewhere.
The assumed equality makes it the same positive value $d_\psi$ at all
frequencies. Therefore $S\Ww_\psi=d_\psi\Id$.
For unnormalized Euclidean sums, $S$ is the adjoint of $\Ww_\psi$, so
$\sum_{b,a}|\Ww_\psi f(b,a)|^2=d_\psi\sum_x|f(x)|^2$.
For our normalized group inner products, the adjoint is instead
$S/(p-1)$ and
\[
 \norm{\Ww_\psi f}_{\CC[G_p]}^2
     =\frac{d_\psi}{p-1}\norm f_{\CC[\FF_p]}^2.
\]
Accordingly the isometry factors multiplying $\Ww_\psi$ are
$1/\sqrt{d_\psi}$ and $\sqrt{(p-1)/d_\psi}$, respectively. Inversion
is a left inverse, or an inverse restricted to the transform's range.
\item Program~\ref{wavelet-transform-listing-finite-body} uses one
FFT convolution per dilation, for $O(p^2\log p)$ operations and
$O(p^2)$ output storage. The literal triple sum would cost $O(p^3)$.
Program~\ref{listing-transfo-inverse-finite-body} performs synthesis in
frequency coordinates and divides by the frame constant. Its default
mode recovers the mean-zero component; the full-space mode checks the
tight-frame condition. Direct complex-valued inner products and independent
reconstruction tests verify the conjugations and scaling.
\begin{listing}[H]
\PythonCode{wavelet-transform-listing-finite-body}
\caption{Affine-group wavelet analysis by ordinary FFTs.}
\label{wavelet-transform-listing-finite-body}
\end{listing}
\begin{listing}[H]
\PythonCode{listing-transfo-inverse-finite-body}
\caption{Wavelet reconstruction with explicit mean-zero and full-space conventions.}
\label{listing-transfo-inverse-finite-body}
\end{listing}
\end{enumerate}
\end{correction}
